\chapter{实验结论与个人反思}

\section{核心实验结论}

本章对应课程评分标准中的“实验结论与个人反思”。本文完成了合规数据集选择、数据预处理、EDA 可视化、分类、回归和聚类三类任务。分类任务得到一个偏召回的高风险筛查方案；回归任务显示数字依赖得分可以被当前数字行为特征较好预测；生产力得分预测较弱，作为负结果分析；聚类任务形成了三类探索性数字生活方式画像。

\section{数据规律与应用价值}

从已有图表和模型结果看，高设备使用、高通知、高社交媒体使用与高风险标记和数字依赖得分存在一定关联。数字依赖预测效果好，说明设备时长、通知数和解锁次数等行为变量对该目标有较强表征能力。生产力得分预测较弱，说明当前特征不足以稳定解释该目标。聚类可以形成“高社交媒体使用型”“高设备依赖型”和“低负荷均衡型”三类画像，但边界不强。

这些结果可为数字生活方式管理提供参考，例如控制设备使用时长、管理通知干扰、合理安排社交媒体使用、保证睡眠和身体活动，以及对高风险初筛结果进行人工复核。

\section{模型优点、不足与局限性}

分类模型通过阈值调优提升了高风险召回，但误报也随之增加。回归模型对数字依赖得分表现较好，但对生产力得分表现较弱。聚类画像具有解释性，但 Silhouette=0.1860，说明聚类边界较弱。数据层面，由于数据集为合成数据，本文结果不能用于医学诊断或现实因果结论。

\section{后续改进方向}

后续可以从三方面改进。第一，使用真实、合规、匿名化的数据验证模型稳定性。第二，引入更细粒度的时间特征，例如夜间设备使用、通知集中度和连续睡眠变化。第三，进一步比较模型解释方法、阈值选择策略和人工复核流程，使筛查结果更适合实际管理场景。

\section{课程学习反思}

通过本项目，我进一步认识到数据合规性、目标泄漏控制和指标选择的重要性。分类任务不能只看 Accuracy，而要结合业务目标关注 Recall 和 F1；回归任务不能只挑效果好的目标，也要诚实解释弱预测和负结果；聚类任务不能因为能分成三类就写成天然群体边界。报告复现性、代码组织、结果文件命名和图表管理同样会影响最终质量。
